%Una explicaci\'on breve del contenido del informe, es decir, detalla: Prop\'osito, Estructura del Documento, Descripci\'on (muy breve) del Problema y Motivaci\'on.

En este informe se realizará una presentación de los trabajos más importantes respecto al problema de Truck Drone Parcel Delivery Problem (TDPDP), el cual consiste en encontrar una ruta para un camión que haga entregas a ciertas ubicaciones en el menor tiempo posible, teniendo el apoyo de un dron que permite hacer cargas en paralelo o con menor tiempo respecto al camión.
La motivación de este trabajo nace de análisis como \cite{moshrefjavadi2021applications} donde se explica el gran auge que ha tenido la utilización de drones en el último tiempo, especialmente con motivos de entrega de productos, dado su alto valor económico para las empresas.
A continuación, se presentarán los principales avances desde su concepción inicial en el trabajo de Murray et al. (2014) \cite{Murray2015Flying}, pasando también por variantes definidas en otros trabajos, y comparando de qué manera se resolvían estos problemas en computación.
